\documentclass[12pt]{article}

\usepackage{sbc-template}
\usepackage{graphicx,url}
\usepackage[utf8]{inputenc}
\usepackage[brazil]{babel}
\usepackage{booktabs}

% Macros para dados numéricos (ajustáveis em um único lugar)
\newcommand{\TotalTasks}{594}
\newcommand{\ScoreTestComp}{57}
\newcommand{\CountTrue}{264}
\newcommand{\CountUnknown}{271}
\newcommand{\CountFalse}{56}
\newcommand{\CountFalseFree}{1}
\newcommand{\CountTimeout}{2}
\newcommand{\PctTrue}{44,4}
\newcommand{\PctUnknown}{45,6}
\newcommand{\PctFalse}{9,4}
\newcommand{\CountBugsFixed}{3}
\newcommand{\CountCommits}{35}
\newcommand{\CountFiles}{84}
\newcommand{\LinesAdded}{5.785}
\newcommand{\LinesRemoved}{657}
\newcommand{\UnitTests}{7/7}
\newcommand{\PassTests}{9/9}
\newcommand{\CountPasses}{9}
\newcommand{\LLVMVersion}{16}
\newcommand{\KLEEVersion}{3.x}
\newcommand{\CppStandard}{C++17}
\newcommand{\DockerBase}{Ubuntu~22.04}
\newcommand{\TimeoutVal}{300}

\sloppy

\title{Map2Check 2026: Modernização Sustentável de uma Ferramenta de Verificação de Memory Safety com LLVM~16 e New Pass Manager}

\author{Guilherme L. P. Bernardo\inst{1}, Inácio Viana\inst{1}, Francisco Nobre\inst{2}, Herbert Rocha\inst{1}}

\address{Departamento de Informática e Matemática Aplicada -- Universidade Federal do Rio Grande do Norte (UFRN)\\
  Natal -- RN -- Brasil
\nextinstitute
  Departamento de Engenharia Elétrica -- Universidade Federal de Roraima (UFRR)\\
  Boa Vista -- RR -- Brasil
  \email{bguilherme51@gmail.com, inacioviana.iv@gmail.com,}\\[-2pt]
  {\footnotesize\texttt{\{diego.nobre,herbert.rocha\}@ufrr.br}}
}

\begin{document}

\maketitle

\begin{abstract}
  Map2Check is a hybrid verification tool for C programs that combines symbolic
  execution and fuzzing to detect memory safety vulnerabilities. After a period
  of stagnation since 2019 (LLVM~6.0), the tool underwent a complete stack
  modernization in 2026: migration to LLVM~\LLVMVersion, adoption of the New Pass Manager,
  \CppStandard{} standard, CI/CD infrastructure, static analysis, and sanitizers. This
  paper presents the engineering modernization, reports results from TestComp
  2026 (\TotalTasks{} tasks, score~\ScoreTestComp), and discusses ongoing WebAssembly support.
\end{abstract}

\begin{resumo}
  O Map2Check é uma ferramenta de verificação híbrida para programas em C que
  combina execução simbólica e fuzzing para detectar vulnerabilidades de
  memory safety. Após um período de estagnação desde 2019 (LLVM~6.0), a
  ferramenta passou por uma modernização completa da stack em 2026: migração
  para LLVM~\LLVMVersion, adoção do New Pass Manager, padrão \CppStandard, infraestrutura de
  CI/CD, análise estática e sanitizers. Este artigo apresenta a modernização
  de engenharia, reporta resultados do TestComp~2026 (\TotalTasks{} tarefas, score~\ScoreTestComp)
  e discute o suporte em desenvolvimento a WebAssembly.
\end{resumo}


\section{Introdução e Motivação}
\label{sec:intro}

\textbf{[PARÁGRAFO -- Problema de memory safety em C/C++]}
Abrir com estatísticas de vulnerabilidades de memória (CWEs 119, 416, 415,
401, 787). Memory-safety bugs continuam entre as principais causas de CVEs
exploráveis~
.

\textbf{[PARÁGRAFO -- Verificação híbrida como abordagem]}
Breve explicação da combinação análise estática + dinâmica (fuzzing +
execução simbólica). Citar trabalhos de ponta: FuSeBMC (smart seeds),
Symbiotic (slicing + KLEE)~
.

\textbf{[PARÁGRAFO -- Posicionamento do Map2Check]}
Map2Check como verificador híbrido consolidado. Menção às publicações
anteriores (TACAS 2016, 2018, 2020)~
.
O problema: estagnação desde 2019 (LLVM~6.0, Ubuntu~16.04, último release
nov/2019).

\textbf{[PARÁGRAFO -- Contribuições desta versão]}
Lista de contribuições: (1) migração LLVM~6$\rightarrow$\LLVMVersion, New Pass
Manager, \CppStandard; (2) infraestrutura CI/CD, Docker, sanitizers, análise
estática; (3) validação em benchmarks reais: TestComp~2026 (Heap, \TotalTasks{}
tarefas, score~\ScoreTestComp); (4) fortalecimento em cibersegurança (mapeamento CWE);
(5) suporte a WebAssembly em desenvolvimento.


\section{História e Evolução}
\label{sec:historia}

\textbf{[PARÁGRAFO -- Linha do tempo 2016$\rightarrow$2026]}
2016 (TACAS): ``Hunting Memory Bugs'' --- BMC com CBLEC.
2018 (TACAS): migração para LLVM + KLEE~
.
2020 (TACAS): abordagem híbrida --- fuzzing + execução simbólica +
invariantes indutivos com LibFuzzer/KLEE/Crab-LLVM~
.
2019--2025: estagnação.
2026: renascimento.
Figura~\ref{fig:timeline}.

\textbf{[PARÁGRAFO -- O que motivou o renascimento]}
Obsolescência técnica: LLVM~6.0 EOL, Ubuntu~16.04 EOL, KLEE~2.0 sem
suporte. Impossibilidade de competir em SV-COMP/TestComp com stack legada.
Necessidade de sustentabilidade de engenharia.

\begin{figure}[ht]
\centering
\includegraphics[width=.9\textwidth]{img/fig1-timeline}
\caption{Evolução do Map2Check: de 2016 a 2026.}
\label{fig:timeline}
\end{figure}


\section{Arquitetura e Funcionalidades}
\label{sec:arquitetura}

\subsection{Pipeline de Verificação}
\label{subsec:pipeline}

\textbf{[PARÁGRAFO -- Visão geral]}
Fluxo completo: Código C $\rightarrow$ clang-16 $\rightarrow$ LLVM IR
$\rightarrow$ opt (passes) $\rightarrow$ llvm-link $\rightarrow$ runtime
$\rightarrow$ KLEE / LibFuzzer $\rightarrow$ Veredito + Witness.
Figura~\ref{fig:pipeline}.

\begin{figure}[ht]
\centering
\includegraphics[width=.95\textwidth]{img/fig2-pipeline}
\caption{Pipeline de verificação do Map2Check 2026.}
\label{fig:pipeline}
\end{figure}

\subsection{Passes de Instrumentação}
\label{subsec:passes}

\textbf{[PARÁGRAFO -- New Pass Manager]}
Descrição dos 9 passes migrados: AssertPass, TargetPass,
LoopPredAssumePass, Map2CheckLibrary, NonDetPass, TrackBasicBlockPass,
OverflowPass, GenerateAutomataTruePass, MemoryTrackPass.
Destacar: \texttt{static bool isRequired() \{ return true; \}} como fix
crítico~
.
Figura~\ref{fig:passes}.

\begin{figure}[ht]
\centering
\includegraphics[width=.95\textwidth]{img/fig3-passes}
\caption{Os 9 passes de instrumentação do Map2Check.}
\label{fig:passes}
\end{figure}

\subsection{Modernização de Engenharia}
\label{subsec:engenharia}

\textbf{[PARÁGRAFO -- Infraestrutura]}
LLVM~16 + New Pass Manager; C++17 (\texttt{std::filesystem}, structured
bindings); CI/CD (GitHub Actions); Docker (\texttt{Dockerfile.dev},
Ubuntu~22.04); análise estática (clang-tidy, cppcheck); sanitizers
(ASAN, UBSAN, TSAN); OpenSSF Best Practices (em andamento).
Esta infraestrutura pontua nos critérios CTA~
.
Figura~\ref{fig:cicd}.

\begin{figure}[ht]
\centering
\includegraphics[width=.95\textwidth]{img/fig5-cicd}
\caption{Workflow CI/CD e qualidade de artefato.}
\label{fig:cicd}
\end{figure}


\section{Features de Cibersegurança}
\label{sec:ciberseguranca}

\subsection{Mapeamento de Memory Safety para CWEs}
\label{subsec:cwe}

\textbf{[PARÁGRAFO -- Mapeamento]}
CWE-401 (Memory Leak), CWE-416 (Use-After-Free), CWE-415 (Double-Free),
CWE-119/787 (Buffer Overflow). Tabela~\ref{tab:cwe}.

\begin{table}[ht]
\centering
\caption{Mapeamento propriedade SV-COMP $\rightarrow$ CWE $\rightarrow$ pass.}
\label{tab:cwe}
\begin{tabular}{@{}lll@{}}
\toprule
Propriedade SV-COMP & CWE & Pass responsável \\
\midrule
\texttt{valid-free} & CWE-415 & MemoryTrackPass \\
\texttt{valid-deref} & CWE-119/787 & MemoryTrackPass + OverflowPass \\
\texttt{valid-memtrack} & CWE-401 & MemoryTrackPass \\
\texttt{no-overflow} & CWE-190 & OverflowPass \\
\bottomrule
\end{tabular}
\end{table}

\subsection{TestComp 2026 --- Execução Heap}
\label{subsec:testcomp}

\textbf{[PARÁGRAFO -- Contexto e resultados]}
Sub-categoria \texttt{C.coverage-error-call.Heap}. Configuração:
timeout \TimeoutVal{}s, Z3, target \texttt{reach\_error}, Docker single-container.
Resultados em \TotalTasks{} tarefas~
:
Tabela~\ref{tab:testcomp}.
Figura~\ref{fig:testcomp}.

\begin{table}[ht]
\centering
\caption{Resultados TestComp 2026 --- C.coverage-error-call.Heap.}
\label{tab:testcomp}
\begin{tabular}{@{}lrr@{}}
\toprule
Resultado & Contagem & \% \\
\midrule
TRUE & \CountTrue & \PctTrue \% \\
UNKNOWN & \CountUnknown & \PctUnknown \% \\
FALSE & \CountFalse & \PctFalse \% \\
FALSE\_FREE & \CountFalseFree & 0,2 \% \\
TIMEOUT & \CountTimeout & 0,3 \% \\
\midrule
\textbf{Total} & \textbf{\TotalTasks} & \textbf{100 \%} \\
\textbf{Score} & \textbf{\ScoreTestComp} & --- \\
\bottomrule
\end{tabular}
\end{table}

\begin{figure}[ht]
\centering
\includegraphics[width=.7\textwidth]{img/fig4-testcomp2026}
\caption{Resultados consolidados TestComp 2026 --- Heap.}
\label{fig:testcomp}
\end{figure}

\subsection{Smoke Test e Validação de Vereditos}
\label{subsec:smoke}

\textbf{[PARÁGRAFO -- E2E]}
Execução manual em \texttt{loops/array-1.c} (TRUE) e
\texttt{loops/array-2.c} (FALSE + counter-example). Framework de
checkpoint automatizado.

\subsection{Bugs Críticos Encontrados e Corrigidos}
\label{subsec:bugs}

\textbf{[PARÁGRAFO -- Tabela dos 3 bugs]}
Tabela~\ref{tab:bugs}.

\begin{table}[ht]
\centering
\caption{Bugs críticos descobertos durante a migração.}
\label{tab:bugs}
\begin{tabular}{@{}llll@{}}
\toprule
Bug & Sintoma & Causa root & Commit \\
\midrule
KLEE 3.x flags & Arg. desconhecida & Flags removidas na v3.x & \texttt{284aef1} \\
\texttt{isRequired()} & Passes pulados & \texttt{optnone} + New PM & \texttt{0deb84e} \\
Target function & Função errada & env var vs \texttt{cl::opt} & \texttt{0deb84e} \\
\bottomrule
\end{tabular}
\end{table}

\subsection{Suporte a WebAssembly (WASM) --- Em Desenvolvimento}
\label{subsec:wasm}

\textbf{[PARÁGRAFO -- WASM]}
Extensão do pipeline LLVM~16 para análise de módulos WASM. Backend
LLVM-WASM: compilação C $\rightarrow$ LLVM IR $\rightarrow$ WASM
$\rightarrow$ instrumentação $\rightarrow$ verificação. Vetor de ataque:
memory safety em runtimes WASM (Wasmtime, Wasmer). Testes preliminares
em benchmarks SV-COMP adaptados. Status: implementação em andamento.


\section{Demonstração Planejada}
\label{sec:demo}

\textbf{[PARÁGRAFO -- O que será demonstrado]}
Pipeline completo em Docker: compilação $\rightarrow$ instrumentação
$\rightarrow$ execução $\rightarrow$ veredito. Exemplo interativo:
\texttt{loops/array-2.c} mostrando veredito FALSE + counter-example.
Execução do framework TestComp~2026 (subset).

\textbf{[PARÁGRAFO -- Equipamentos]}
Laptop com Docker instalado. Imagem \texttt{map2check-dev}
pré-construída. Subconjunto de benchmarks SV-COMP. Acesso à internet.

\textbf{[PARÁGRAFO -- Vídeo técnico]}
Link do vídeo (5--7 min): instalação, build, execução E2E, resultados
TestComp~2026.


\section{Avaliação, Perspectivas e Conclusão}
\label{sec:conclusao}

\textbf{[PARÁGRAFO -- Resumo]}
Modernização: \CountCommits{} commits, \CountFiles{} arquivos, +\LinesAdded{}/$-$\LinesRemoved{} linhas. \CountPasses{} passes
migrados, \UnitTests{} unit tests, \PassTests{} passes carregando. TestComp~2026 Heap:
\TotalTasks{} tarefas, score~\ScoreTestComp, \CountFalse{} bugs reais.

\textbf{[PARÁGRAFO -- Limitações]}
Taxa de UNKNOWN alta (\PctUnknown\%). Ausência de testes E2E no CI.

\textbf{[PARÁGRAFO -- Perspectivas]}
Aprofundamento do suporte a WebAssembly. Redução da taxa de UNKNOWN.
Testes de integração E2E no CI.

\textbf{[PARÁGRAFO -- Conclusão]}
O Map2Check renasceu como plataforma sustentável de verificação de
memory safety. A modernização de engenharia é pré-requisito para
qualquer avanço futuro.


\section*{Agradecimentos}

Agradecemos ao CNPq e à CAPES pelo apoio financeiro.

\bibliographystyle{sbc}
\bibliography{references}

\newpage
\appendix

\section{Reprodutibilidade}
\label{ap:repro}

\subsection{Build com Docker}

\begin{verbatim}
docker build -f Dockerfile.dev -t map2check-dev .
docker run -it --rm map2check-dev
\end{verbatim}

\subsection{Execução de Benchmarks}

\begin{verbatim}
map2check --property-file coverage-error-call.prp \
  --target-function-name reach_error loops/array-1.c

cd test-comp2026/simulation/
./run_checkpoint.sh --category Heap --timeout \TimeoutVal
./verify_verdicts.sh
\end{verbatim}

\subsection{Declaração de Uso de IA Generativa}
\label{ap:ia}

Ferramentas de Inteligência Artificial Generativa foram utilizadas
como assistentes de escrita e revisão no processo de redação deste
artigo, conforme o Código de Conduta da SBC.

\end{document}
